% OB-SECTION {"id":"related-work","label":"sec-related","level":2,"order":2,"title":"Related Work"}

\section{Related Work}
\label{sec-related}
% OB-EDIT-BEGIN openbarry/markdown@1 id=related-work
Prior work has established that language model outputs often exhibit biases towards specific demographic groups~\cite{santurkar2023}, genders~\cite{delarco2024}, countries~\cite{durmus2023}, or personas~\cite{giorgi2024}.
This homogenization can limit LLM utility on downstream tasks~\cite{agarwal2024,song2024}.
Research into diversifying language model outputs has seen several promising directions, including adjusting decoding methods such as increasing temperature~\cite{shurofry2024,peeperkorn2024}, or modifying training objectives~\cite{li2016diversity,zhang2024diffuse,li2024preserving,lanchantin2025}.

The idea that we should evaluate language models not only on the quality of their generations but also on their diversity is not entirely new.
The 2020 Duolingo Shared Task involved generating several diverse but still correct translations of an input sentence~\cite{mayhew2020}.
Cheng et al.~\cite{cheng2024} highlight the probabilistic nature of commonsense knowledge and proposed an evaluation method that relies on open-ended generation instead of multiple-choice questions.
Numerous studies of NLG systems have proposed metrics for evaluating diversity, based on perplexity~\cite{lewis2017,fan2018}, $n$-gram overlaps~\cite{zhu2018,shu2019,dusek2020}, and embedding distances~\cite{du2019}.
However, existing metrics often correlate poorly with human judgments of content diversity~\cite{tevet2021}.

\obtodo[id=nb-related-metrics, created=2024-11-03T13:40:00.000Z, updated=2025-03-02T10:18:44.000Z, completed=2025-03-02T10:18:44.000Z]{Position NoveltyBench against n-gram, embedding, and entropy diversity metrics.}

In contrast to prior work, \textsc{NoveltyBench} uniquely evaluates an LLM's capacity to generate a set of outputs that are simultaneously diverse and high-quality.
% OB-EDIT-END id=related-work
